\documentclass[11pt]{article}
\newcommand{\ListaTitulo}{Gabarito --- Lista 03}
% Preâmbulo compartilhado — MATD48, Listas de Exercícios 2026
% Usado por Lista{NN}.tex e Gabarito{NN}.tex via % Preâmbulo compartilhado — MATD48, Listas de Exercícios 2026
% Usado por Lista{NN}.tex e Gabarito{NN}.tex via % Preâmbulo compartilhado — MATD48, Listas de Exercícios 2026
% Usado por Lista{NN}.tex e Gabarito{NN}.tex via \input{preamble}

\usepackage[utf8]{inputenc}
\usepackage[T1]{fontenc}
\usepackage[brazil]{babel}
\usepackage{fullpage}
\usepackage{amsmath,amssymb}
\usepackage{enumitem}
\usepackage{xcolor}
\usepackage{fancyhdr}
\usepackage{titlesec}
\usepackage{listings}
\usepackage{hyperref}
\usepackage{booktabs}
\usepackage{graphicx}

\definecolor{matd48azul}{HTML}{0B4F6C}
\definecolor{matd48verde}{HTML}{2E8B57}

\titleformat{\section}{\normalfont\Large\bfseries\color{matd48azul}}{\thesection}{1em}{}
\titleformat{\subsection}{\normalfont\large\bfseries\color{matd48azul}}{\thesubsection}{1em}{}

\providecommand{\ListaTitulo}{Lista de Exercícios}

\setlength{\headheight}{14pt}
\pagestyle{fancy}
\fancyhf{}
\lhead{\small MATD48 --- Planejamento de Experimentos A}
\rhead{\small \ListaTitulo}
\cfoot{\thepage}
\renewcommand{\headrulewidth}{0.4pt}

\lstset{
  basicstyle=\ttfamily\small,
  breaklines=true,
  frame=single,
  columns=fullflexible,
  backgroundcolor=\color{gray!8},
  commentstyle=\color{matd48verde},
  keywordstyle=\color{matd48azul}\bfseries,
  language=R,
  extendedchars=true,
  literate=%
    {á}{{\'a}}1 {à}{{\`a}}1 {â}{{\^a}}1 {ã}{{\~a}}1
    {é}{{\'e}}1 {ê}{{\^e}}1 {í}{{\'i}}1 {ó}{{\'o}}1
    {ô}{{\^o}}1 {õ}{{\~o}}1 {ú}{{\'u}}1 {ç}{{\c c}}1
    {Á}{{\'A}}1 {À}{{\`A}}1 {Â}{{\^A}}1 {Ã}{{\~A}}1
    {É}{{\'E}}1 {Ê}{{\^E}}1 {Í}{{\'I}}1 {Ó}{{\'O}}1
    {Ô}{{\^O}}1 {Õ}{{\~O}}1 {Ú}{{\'U}}1 {Ç}{{\c C}}1
}

\hypersetup{colorlinks=true, linkcolor=matd48azul, urlcolor=matd48azul, citecolor=matd48azul}

\newcommand{\cabecalho}[2]{%
  \begin{center}
    {\Large\bfseries\color{matd48azul} #1}\\[0.3em]
    {\large #2}\\[0.3em]
    \rule{\linewidth}{0.6pt}
  \end{center}
  \vspace{0.5em}
}

\newenvironment{questao}[1]{\vspace{1em}\noindent\textbf{Questão #1.}\hspace{0.4em}}{\par}
\newenvironment{solucao}{\vspace{0.3em}\noindent\textit{Solução.}\hspace{0.4em}\color{black!85}}{\par\vspace{0.5em}}


\usepackage[utf8]{inputenc}
\usepackage[T1]{fontenc}
\usepackage[brazil]{babel}
\usepackage{fullpage}
\usepackage{amsmath,amssymb}
\usepackage{enumitem}
\usepackage{xcolor}
\usepackage{fancyhdr}
\usepackage{titlesec}
\usepackage{listings}
\usepackage{hyperref}
\usepackage{booktabs}
\usepackage{graphicx}

\definecolor{matd48azul}{HTML}{0B4F6C}
\definecolor{matd48verde}{HTML}{2E8B57}

\titleformat{\section}{\normalfont\Large\bfseries\color{matd48azul}}{\thesection}{1em}{}
\titleformat{\subsection}{\normalfont\large\bfseries\color{matd48azul}}{\thesubsection}{1em}{}

\providecommand{\ListaTitulo}{Lista de Exercícios}

\setlength{\headheight}{14pt}
\pagestyle{fancy}
\fancyhf{}
\lhead{\small MATD48 --- Planejamento de Experimentos A}
\rhead{\small \ListaTitulo}
\cfoot{\thepage}
\renewcommand{\headrulewidth}{0.4pt}

\lstset{
  basicstyle=\ttfamily\small,
  breaklines=true,
  frame=single,
  columns=fullflexible,
  backgroundcolor=\color{gray!8},
  commentstyle=\color{matd48verde},
  keywordstyle=\color{matd48azul}\bfseries,
  language=R,
  extendedchars=true,
  literate=%
    {á}{{\'a}}1 {à}{{\`a}}1 {â}{{\^a}}1 {ã}{{\~a}}1
    {é}{{\'e}}1 {ê}{{\^e}}1 {í}{{\'i}}1 {ó}{{\'o}}1
    {ô}{{\^o}}1 {õ}{{\~o}}1 {ú}{{\'u}}1 {ç}{{\c c}}1
    {Á}{{\'A}}1 {À}{{\`A}}1 {Â}{{\^A}}1 {Ã}{{\~A}}1
    {É}{{\'E}}1 {Ê}{{\^E}}1 {Í}{{\'I}}1 {Ó}{{\'O}}1
    {Ô}{{\^O}}1 {Õ}{{\~O}}1 {Ú}{{\'U}}1 {Ç}{{\c C}}1
}

\hypersetup{colorlinks=true, linkcolor=matd48azul, urlcolor=matd48azul, citecolor=matd48azul}

\newcommand{\cabecalho}[2]{%
  \begin{center}
    {\Large\bfseries\color{matd48azul} #1}\\[0.3em]
    {\large #2}\\[0.3em]
    \rule{\linewidth}{0.6pt}
  \end{center}
  \vspace{0.5em}
}

\newenvironment{questao}[1]{\vspace{1em}\noindent\textbf{Questão #1.}\hspace{0.4em}}{\par}
\newenvironment{solucao}{\vspace{0.3em}\noindent\textit{Solução.}\hspace{0.4em}\color{black!85}}{\par\vspace{0.5em}}


\usepackage[utf8]{inputenc}
\usepackage[T1]{fontenc}
\usepackage[brazil]{babel}
\usepackage{fullpage}
\usepackage{amsmath,amssymb}
\usepackage{enumitem}
\usepackage{xcolor}
\usepackage{fancyhdr}
\usepackage{titlesec}
\usepackage{listings}
\usepackage{hyperref}
\usepackage{booktabs}
\usepackage{graphicx}

\definecolor{matd48azul}{HTML}{0B4F6C}
\definecolor{matd48verde}{HTML}{2E8B57}

\titleformat{\section}{\normalfont\Large\bfseries\color{matd48azul}}{\thesection}{1em}{}
\titleformat{\subsection}{\normalfont\large\bfseries\color{matd48azul}}{\thesubsection}{1em}{}

\providecommand{\ListaTitulo}{Lista de Exercícios}

\setlength{\headheight}{14pt}
\pagestyle{fancy}
\fancyhf{}
\lhead{\small MATD48 --- Planejamento de Experimentos A}
\rhead{\small \ListaTitulo}
\cfoot{\thepage}
\renewcommand{\headrulewidth}{0.4pt}

\lstset{
  basicstyle=\ttfamily\small,
  breaklines=true,
  frame=single,
  columns=fullflexible,
  backgroundcolor=\color{gray!8},
  commentstyle=\color{matd48verde},
  keywordstyle=\color{matd48azul}\bfseries,
  language=R,
  extendedchars=true,
  literate=%
    {á}{{\'a}}1 {à}{{\`a}}1 {â}{{\^a}}1 {ã}{{\~a}}1
    {é}{{\'e}}1 {ê}{{\^e}}1 {í}{{\'i}}1 {ó}{{\'o}}1
    {ô}{{\^o}}1 {õ}{{\~o}}1 {ú}{{\'u}}1 {ç}{{\c c}}1
    {Á}{{\'A}}1 {À}{{\`A}}1 {Â}{{\^A}}1 {Ã}{{\~A}}1
    {É}{{\'E}}1 {Ê}{{\^E}}1 {Í}{{\'I}}1 {Ó}{{\'O}}1
    {Ô}{{\^O}}1 {Õ}{{\~O}}1 {Ú}{{\'U}}1 {Ç}{{\c C}}1
}

\hypersetup{colorlinks=true, linkcolor=matd48azul, urlcolor=matd48azul, citecolor=matd48azul}

\newcommand{\cabecalho}[2]{%
  \begin{center}
    {\Large\bfseries\color{matd48azul} #1}\\[0.3em]
    {\large #2}\\[0.3em]
    \rule{\linewidth}{0.6pt}
  \end{center}
  \vspace{0.5em}
}

\newenvironment{questao}[1]{\vspace{1em}\noindent\textbf{Questão #1.}\hspace{0.4em}}{\par}
\newenvironment{solucao}{\vspace{0.3em}\noindent\textit{Solução.}\hspace{0.4em}\color{black!85}}{\par\vspace{0.5em}}


\begin{document}

\cabecalho{Gabarito --- Lista de Exercícios 03}{Gauss-Markov, matriz de projeção, modelo particionado e formas quadráticas (Capítulo 2 / Aula 03)}

\begin{questao}{1} \textbf{(Dedução --- propriedades da matriz de projeção)}
\end{questao}
\begin{solucao}
\begin{enumerate}[label=(\alph*)]
  \item $\mathbf{P}_X' = \big(\mathbf{X}(\mathbf{X}'\mathbf{X})^{-1}\mathbf{X}'\big)' =
    (\mathbf{X}')'\big((\mathbf{X}'\mathbf{X})^{-1}\big)'\mathbf{X}' =
    \mathbf{X}\big((\mathbf{X}'\mathbf{X})^{-1}\big)'\mathbf{X}'$. Como $\mathbf{X}'\mathbf{X}$ é
    simétrica, sua inversa também é simétrica, logo $\big((\mathbf{X}'\mathbf{X})^{-1}\big)' =
    (\mathbf{X}'\mathbf{X})^{-1}$, e $\mathbf{P}_X' = \mathbf{X}(\mathbf{X}'\mathbf{X})^{-1}\mathbf{X}' =
    \mathbf{P}_X$.
  \item
  $$
  \mathbf{P}_X\mathbf{P}_X = \mathbf{X}(\mathbf{X}'\mathbf{X})^{-1}\mathbf{X}'\,\mathbf{X}(\mathbf{X}'\mathbf{X})^{-1}\mathbf{X}'
  = \mathbf{X}(\mathbf{X}'\mathbf{X})^{-1}\big[(\mathbf{X}'\mathbf{X})(\mathbf{X}'\mathbf{X})^{-1}\big]\mathbf{X}'
  = \mathbf{X}(\mathbf{X}'\mathbf{X})^{-1}\mathbf{I}_p\,\mathbf{X}' = \mathbf{P}_X.
  $$
  O termo entre colchetes é $\mathbf{I}_p$ porque é uma matriz multiplicada por sua própria
  inversa.
  \item Simetria: $(\mathbf{I}_n-\mathbf{P}_X)' = \mathbf{I}_n' - \mathbf{P}_X' = \mathbf{I}_n -
    \mathbf{P}_X$ (usando o item (a)). Idempotência:
    $(\mathbf{I}_n-\mathbf{P}_X)(\mathbf{I}_n-\mathbf{P}_X) = \mathbf{I}_n - \mathbf{P}_X -
    \mathbf{P}_X + \mathbf{P}_X\mathbf{P}_X = \mathbf{I}_n - 2\mathbf{P}_X + \mathbf{P}_X =
    \mathbf{I}_n - \mathbf{P}_X$ (usando o item (b) na última passagem).
  \item Se $\mathbf{P}_X\mathbf{v} = \lambda\mathbf{v}$ com $\mathbf{v}\neq\mathbf{0}$, então
    $\mathbf{P}_X\mathbf{P}_X\mathbf{v} = \mathbf{P}_X(\lambda\mathbf{v}) = \lambda^2\mathbf{v}$.
    Mas, por idempotência (item (b)), $\mathbf{P}_X\mathbf{P}_X\mathbf{v} = \mathbf{P}_X\mathbf{v}
    = \lambda\mathbf{v}$. Logo $\lambda^2\mathbf{v} = \lambda\mathbf{v}$, ou seja,
    $\lambda(\lambda-1)\mathbf{v}=\mathbf{0}$; como $\mathbf{v}\neq\mathbf{0}$, segue
    $\lambda(\lambda-1)=0$, isto é, $\lambda \in \{0,1\}$. O traço de uma matriz é a soma de seus
    autovalores (com multiplicidade); como só há autovalores 0 e 1, o traço conta exatamente o
    número de autovalores iguais a 1, que é, por definição, o posto da matriz. Logo
    $\mathrm{tr}(\mathbf{P}_X) = \mathrm{posto}(\mathbf{P}_X)$.
\end{enumerate}
\end{solucao}

\begin{questao}{2} \textbf{(Dedução --- variância de $\hat{\boldsymbol\beta}$)}
\end{questao}
\begin{solucao}
\begin{enumerate}[label=(\alph*)]
  \item $\hat{\boldsymbol\beta} = \mathbf{A}\mathbf{Y}$ com
    $\mathbf{A} = (\mathbf{X}'\mathbf{X})^{-1}\mathbf{X}'$.
  \item
  $$
  \mathrm{Cov}(\hat{\boldsymbol\beta}) = \mathbf{A}\,\mathrm{Cov}(\mathbf{Y})\,\mathbf{A}' =
  (\mathbf{X}'\mathbf{X})^{-1}\mathbf{X}'\ \big(\sigma^2\mathbf{I}_n\big)\
  \mathbf{X}(\mathbf{X}'\mathbf{X})^{-1} = \sigma^2 (\mathbf{X}'\mathbf{X})^{-1}
  \underbrace{\mathbf{X}'\mathbf{X}}_{}(\mathbf{X}'\mathbf{X})^{-1}
  $$
  $$
  = \sigma^2 (\mathbf{X}'\mathbf{X})^{-1}.
  $$
  (Usamos $\mathbf{A}' = \mathbf{X}\big((\mathbf{X}'\mathbf{X})^{-1}\big)' =
  \mathbf{X}(\mathbf{X}'\mathbf{X})^{-1}$, pois $(\mathbf{X}'\mathbf{X})^{-1}$ é simétrica, e o
  produto do meio $(\mathbf{X}'\mathbf{X})^{-1}(\mathbf{X}'\mathbf{X})(\mathbf{X}'\mathbf{X})^{-1}$
  se reduz a $(\mathbf{X}'\mathbf{X})^{-1}$.)
  \item $\mathrm{Var}(\hat\beta_j) = \sigma^2\, [(\mathbf{X}'\mathbf{X})^{-1}]_{jj}$, o $j$-ésimo
    elemento da diagonal de $(\mathbf{X}'\mathbf{X})^{-1}$, multiplicado por $\sigma^2$.
\end{enumerate}
\end{solucao}

\begin{questao}{3} \textbf{(Psicologia --- matriz de projeção numérica)}
\end{questao}
\begin{solucao}
\begin{enumerate}[label=(\alph*)]
  \item
  $$
  \mathbf{X} = \begin{bmatrix} 1 & 1 \\ 1 & 2 \\ 1 & 3 \\ 1 & 4 \end{bmatrix}.
  $$
  \item $\mathbf{X}'\mathbf{X} = \begin{bmatrix} 4 & 10 \\ 10 & 30 \end{bmatrix}$, com
  $\det = 4(30) - 10^2 = 20$, então $(\mathbf{X}'\mathbf{X})^{-1} = \frac{1}{20}\begin{bmatrix} 30
  & -10 \\ -10 & 4 \end{bmatrix} = \begin{bmatrix} 1{,}5 & -0{,}5 \\ -0{,}5 & 0{,}2 \end{bmatrix}$.
  Calculando $\mathbf{P}_X = \mathbf{X}(\mathbf{X}'\mathbf{X})^{-1}\mathbf{X}'$:
  $$
  \mathbf{P}_X = \begin{bmatrix}
  0{,}7 & 0{,}4 & 0{,}1 & -0{,}2 \\
  0{,}4 & 0{,}3 & 0{,}2 & 0{,}1 \\
  0{,}1 & 0{,}2 & 0{,}3 & 0{,}4 \\
  -0{,}2 & 0{,}1 & 0{,}4 & 0{,}7
  \end{bmatrix}.
  $$
  \item Por exemplo, a entrada $(1,1)$ de $\mathbf{P}_X\mathbf{P}_X$:
  $0{,}7^2 + 0{,}4^2 + 0{,}1^2 + (-0{,}2)^2 = 0{,}49+0{,}16+0{,}01+0{,}04 = 0{,}70$, que coincide
  com a entrada $(1,1)$ de $\mathbf{P}_X$ --- consistente com a idempotência provada na Questão 1.
  $\mathrm{tr}(\mathbf{P}_X) = 0{,}7+0{,}3+0{,}3+0{,}7 = 2$, que é exatamente
  $\mathrm{posto}(\mathbf{X}) = 2$.
  \item Alavancagens: $h_{11}=0{,}7$, $h_{22}=0{,}3$, $h_{33}=0{,}3$, $h_{44}=0{,}7$. As
    observações \textbf{1 e 4} (os extremos $x=1$ e $x=4$) têm alavancagem maior. Geometricamente,
    a reta de mínimos quadrados é "presa" com mais força pelos pontos mais distantes do centro
    $\bar x$: mover um ponto extremo desloca a reta ajustada muito mais do que mover um ponto
    central, porque o braço de alavanca (a distância $x_i - \bar x$) é maior.
  \item
  $$
  \hat{\mathbf{Y}} = \mathbf{P}_X\mathbf{Y} = (2{,}7,\ 4{,}9,\ 7{,}1,\ 9{,}3)', \qquad
  \mathrm{SQE} = \mathbf{Y}'(\mathbf{I}_4-\mathbf{P}_X)\mathbf{Y} = \sum_i (y_i-\hat y_i)^2 =
  0{,}3^2+0{,}1^2+(-1{,}1)^2+0{,}7^2 = 1{,}8.
  $$
\end{enumerate}
\end{solucao}

\begin{questao}{4} \textbf{(Ciência de dados --- soma de quadrados extra)}
\end{questao}
\begin{solucao}
\begin{enumerate}[label=(\alph*)]
  \item $5000 - 4200 = \textbf{800}$.
  \item Numerador: $p - p_1 = 3-2 = \textbf{1}$. Denominador: $n-p = 40-3=\textbf{37}$.
  \item
  $$
  F = \frac{800/1}{4200/37} = \frac{800}{113{,}51} \approx \textbf{7{,}05}.
  $$
  Como $7{,}05 > 4{,}11$, \textbf{rejeitamos $H_0$} ao nível de 5\%: a covariável adicional
  contribui significativamente para explicar a resposta, além do que o modelo reduzido já
  explica.
  \item $\mathrm{SQE}(\text{reduzido}) - \mathrm{SQE}(\text{completo}) =
    \mathbf{Y}'(\mathbf{P}_X - \mathbf{P}_{X_1})\mathbf{Y}$ --- a forma quadrática associada à
    diferença entre as duas matrizes de projeção.
\end{enumerate}
\end{solucao}

\begin{questao}{5} \textbf{(Distribuição de formas quadráticas)}
\end{questao}
\begin{solucao}
\begin{enumerate}[label=(\alph*)]
  \item
  $$
  \mathbf{A}\mathbf{A} = \begin{bmatrix} 0{,}5 & 0{,}5 \\ 0{,}5 & 0{,}5\end{bmatrix}
  \begin{bmatrix} 0{,}5 & 0{,}5 \\ 0{,}5 & 0{,}5\end{bmatrix} =
  \begin{bmatrix} 0{,}5 & 0{,}5 \\ 0{,}5 & 0{,}5\end{bmatrix} = \mathbf{A}.
  $$
  $\mathbf{A}$ é idempotente. $\mathrm{tr}(\mathbf{A}) = 0{,}5+0{,}5=1$, e pela Questão 1(d),
  $\mathrm{posto}(\mathbf{A}) = \mathrm{tr}(\mathbf{A}) = \textbf{1}$.
  \item Como $\mathbf{A}\boldsymbol\mu = \mathbf{A}(0,0)' = (0,0)' = \mathbf{0}$, pelo teorema de
    distribuição de formas quadráticas idempotentes, $\mathbf{Y}'\mathbf{A}\mathbf{Y}/\sigma^2
    \sim \chi^2_1$ (\textbf{qui-quadrado central} com 1 grau de liberdade, igual ao posto de
    $\mathbf{A}$).
  \item $\mathbf{A}\boldsymbol\mu = \begin{bmatrix}0{,}5(3)+0{,}5(-3) \\ 0{,}5(3)+0{,}5(-3)\end{bmatrix}
    = \begin{bmatrix}0\\0\end{bmatrix} = \mathbf{0}$. Mesmo com $\boldsymbol\mu \neq \mathbf{0}$,
    a resposta do item (b) \textbf{não muda}: a distribuição continua sendo qui-quadrado
    \textbf{central} com 1 grau de liberdade. O que determina centralidade não é $\boldsymbol\mu$
    ser nulo, mas sim $\mathbf{A}\boldsymbol\mu$ ser nulo --- e isso acontece sempre que
    $\boldsymbol\mu$ está no espaço-nulo de $\mathbf{A}$. Geometricamente, $\mathbf{A}$ projeta
    sobre a direção $(1,1)/\sqrt2$; o vetor $(3,-3)'$ é ortogonal a essa direção (está no
    complemento ortogonal, o espaço-nulo de $\mathbf{A}$), então a projeção de $\boldsymbol\mu$
    por $\mathbf{A}$ se anula, apesar de $\boldsymbol\mu$ em si não ser o vetor zero.
  \item $\mathbf{A}\boldsymbol\mu = \begin{bmatrix}0{,}5(2)+0{,}5(2)\\0{,}5(2)+0{,}5(2)\end{bmatrix}
    = \begin{bmatrix}2\\2\end{bmatrix} = \boldsymbol\mu$ (porque $(2,2)'$ já está na direção
    $(1,1)$, o espaço-coluna de $\mathbf{A}$). $\boldsymbol\mu'\mathbf{A}\boldsymbol\mu =
    (2,2)\cdot(2,2)' = 4+4=8$. Como $\mathbf{A}\boldsymbol\mu\neq\mathbf{0}$, a distribuição agora
    é \textbf{qui-quadrado não central} com 1 grau de liberdade e parâmetro de não centralidade
    $\boldsymbol\mu'\mathbf{A}\boldsymbol\mu/\sigma^2 = 8/\sigma^2$.
\end{enumerate}
\end{solucao}

\begin{questao}{6} \textbf{(Ciência de dados --- exercício computacional em R)}
\end{questao}
\begin{solucao}
\begin{enumerate}[label=(\alph*)]
  \item
\begin{lstlisting}
gl_extra <- qr(X)$rank - qr(X1)$rank
\end{lstlisting}
  \item
\begin{lstlisting}
gl_residuo <- n - qr(X)$rank
\end{lstlisting}
  \item Bastaria adicionar a segunda covariável à fórmula de \texttt{X} (por exemplo,
    \texttt{model.matrix(\textasciitilde{} horas\_sono + cafeina\_mg + temperatura, data =
    dados\_sono)}); o resto do código --- incluindo o cálculo de \texttt{P}, \texttt{sqe\_completo}
    e \texttt{gl\_extra} --- não muda de forma nenhuma, porque \texttt{gl\_extra} já é calculado
    genericamente como a diferença de postos entre os dois modelos (que passaria de 1 para 2
    automaticamente).
\end{enumerate}
\end{solucao}

\begin{questao}{7} \textbf{(Dedução --- Teorema de Frisch-Waugh-Lovell)}
\end{questao}
\begin{solucao}
\begin{enumerate}[label=(\alph*)]
  \item Da primeira linha de blocos, $\mathbf{X}_1'\mathbf{X}_1\hat{\boldsymbol\beta}_1 +
    \mathbf{X}_1'\mathbf{X}_2\hat{\boldsymbol\beta}_2 = \mathbf{X}_1'\mathbf{Y}$, logo
    $$
    \hat{\boldsymbol\beta}_1 = (\mathbf{X}_1'\mathbf{X}_1)^{-1}\big(\mathbf{X}_1'\mathbf{Y} -
    \mathbf{X}_1'\mathbf{X}_2\hat{\boldsymbol\beta}_2\big).
    $$
  \item Substituindo na segunda linha de blocos, $\mathbf{X}_2'\mathbf{X}_1\hat{\boldsymbol\beta}_1
    + \mathbf{X}_2'\mathbf{X}_2\hat{\boldsymbol\beta}_2 = \mathbf{X}_2'\mathbf{Y}$:
    $$
    \mathbf{X}_2'\mathbf{X}_1(\mathbf{X}_1'\mathbf{X}_1)^{-1}\mathbf{X}_1'\mathbf{Y} -
    \mathbf{X}_2'\underbrace{\mathbf{X}_1(\mathbf{X}_1'\mathbf{X}_1)^{-1}\mathbf{X}_1'}_{=\mathbf{P}_{X_1}}\mathbf{X}_2\hat{\boldsymbol\beta}_2
    + \mathbf{X}_2'\mathbf{X}_2\hat{\boldsymbol\beta}_2 = \mathbf{X}_2'\mathbf{Y}.
    $$
    Reagrupando os termos em $\hat{\boldsymbol\beta}_2$ de um lado e passando
    $\mathbf{X}_2'\mathbf{P}_{X_1}\mathbf{Y}$ para o outro:
    $$
    \mathbf{X}_2'(\mathbf{X}_2 - \mathbf{P}_{X_1}\mathbf{X}_2)\hat{\boldsymbol\beta}_2 =
    \mathbf{X}_2'(\mathbf{Y} - \mathbf{P}_{X_1}\mathbf{Y})
    \ \Longrightarrow\
    \mathbf{X}_2'\mathbf{M}_1\mathbf{X}_2\,\hat{\boldsymbol\beta}_2 = \mathbf{X}_2'\mathbf{M}_1\mathbf{Y},
    $$
    usando $\mathbf{M}_1 = \mathbf{I}_n - \mathbf{P}_{X_1}$ nos dois lados.
  \item Como $\mathbf{M}_1' = \mathbf{M}_1$ e $\mathbf{M}_1\mathbf{M}_1 = \mathbf{M}_1$
    (idempotência), temos $\mathbf{M}_1 = \mathbf{M}_1'\mathbf{M}_1$. Logo
    $$
    \mathbf{X}_2'\mathbf{M}_1\mathbf{X}_2 = \mathbf{X}_2'\mathbf{M}_1'\mathbf{M}_1\mathbf{X}_2 =
    (\mathbf{M}_1\mathbf{X}_2)'(\mathbf{M}_1\mathbf{X}_2) = \tilde{\mathbf{X}}_2'\tilde{\mathbf{X}}_2,
    $$
    e, pelo mesmo argumento, $\mathbf{X}_2'\mathbf{M}_1\mathbf{Y} =
    (\mathbf{M}_1\mathbf{X}_2)'(\mathbf{M}_1\mathbf{Y}) = \tilde{\mathbf{X}}_2'\tilde{\mathbf{Y}}$.
    Substituindo no resultado do item (b):
    $$
    \tilde{\mathbf{X}}_2'\tilde{\mathbf{X}}_2\,\hat{\boldsymbol\beta}_2 =
    \tilde{\mathbf{X}}_2'\tilde{\mathbf{Y}}
    \ \Longrightarrow\
    \hat{\boldsymbol\beta}_2 = (\tilde{\mathbf{X}}_2'\tilde{\mathbf{X}}_2)^{-1}\tilde{\mathbf{X}}_2'\tilde{\mathbf{Y}},
    $$
    que é exatamente a fórmula de mínimos quadrados da regressão simples (sem intercepto, já que
    resíduos de uma regressão com intercepto em $\mathbf{X}_1$ já têm média zero) de
    $\tilde{\mathbf{Y}}$ em $\tilde{\mathbf{X}}_2$. \textbf{C.Q.D.}
  \item
\begin{lstlisting}
Y_tilde  <- resid(lm(tempo_reacao ~ horas_sono, data = dados_sono))
X2_tilde <- resid(lm(cafeina_mg  ~ horas_sono, data = dados_sono))
\end{lstlisting}
    Os dois valores do vetor final devem coincidir porque é exatamente isso que os itens
    (a)--(c) provaram algebricamente: $\hat{\boldsymbol\beta}_2$ (aqui, o coeficiente de
    \texttt{cafeina\_mg}) calculado pela regressão múltipla completa é, por construção das
    equações normais, idêntico ao coeficiente da regressão simples entre os resíduos
    $\tilde{\mathbf{Y}}$ e $\tilde{\mathbf{X}}_2$ --- não é uma coincidência numérica, é uma
    identidade algébrica.
  \item Se \texttt{cafeina\_mg} fosse perfeitamente independente de \texttt{horas\_sono} por
    desenho (sorteada sem relação com o sono), regredir \texttt{cafeina\_mg} em
    \texttt{horas\_sono} produziria um coeficiente angular próximo de zero, e os resíduos
    $\tilde{\mathbf{X}}_2$ seriam quase idênticos à variável original $\mathbf{X}_2$ (a
    covariável "sobra" quase intacta, porque \texttt{horas\_sono} não explica nada dela). Isso
    \textbf{não muda} a conclusão de que $\hat\beta_2$ coincide nos dois métodos --- essa
    identidade vale sempre, por construção algébrica, independentemente de haver ou não
    correlação entre $\mathbf{X}_1$ e $\mathbf{X}_2$. O que muda é a \emph{interpretação}
    prática: sob essa independência por desenho, incluir \texttt{horas\_sono} no modelo deixa de
    ser necessário para estimar $\hat\beta_2$ sem viés (a regressão simples de
    \texttt{tempo\_reacao} em \texttt{cafeina\_mg} já daria, em expectativa, o mesmo coeficiente)
    --- exatamente o ponto da Seção de conexão do FWL com planejamento experimental do Capítulo 2
    do livro.
\end{enumerate}
\end{solucao}

\end{document}
